\documentclass[12pt]{amsart}
\usepackage{amsmath, amssymb, amsthm}
\usepackage[margin=2.5cm]{geometry}
\usepackage{xr-hyper}
\usepackage{hyperref}
\usepackage{xurl}
\usepackage{booktabs}
\emergencystretch=2em

\newtheorem{theorem}{Theorem}
\newtheorem{corollary}[theorem]{Corollary}
\newtheorem{proposition}[theorem]{Proposition}
\newtheorem{lemma}[theorem]{Lemma}
\newtheorem{observation}{Observation}
\theoremstyle{remark}
\newtheorem{remark}{Remark}
\renewcommand{\thesection}{S\arabic{section}}
\renewcommand{\thetheorem}{S\arabic{theorem}}
\renewcommand{\theobservation}{S\arabic{observation}}
\renewcommand{\theremark}{S\arabic{remark}}
\renewcommand{\thetable}{S\arabic{table}}
\renewcommand{\theequation}{S\arabic{equation}}

\newcommand{\E}{\mathbb{E}}
\newcommand{\Et}{\widetilde{\mathbb{E}}}
\newcommand{\eps}{\varepsilon}
\newcommand{\T}{\mathbb{T}}
\newcommand{\1}{\mathbf{1}}
\newcommand{\sinc}{\operatorname{sinc}}
\DeclareMathOperator{\Var}{Var}
\externaldocument[M-]{arxiv_small_gap_hump}

\title[Supplementary material]{Supplementary material for\\ ``The hump
between two close zeros: random unitary matrices and the Riemann zeta
function''}

\author{U\u{g}ur Sezen}
\address{Independent researcher, Istanbul, Turkey}
\email{usezen@aof.anadolu.edu.tr}

\date{\today}

\begin{document}

\maketitle

Numbers without the prefix S (theorems, lemmas, equations, sections,
tables) refer to the main text; notation is that of its
Section~\ref{M-sec:setup}.

%======================================================================
\section{Complements to Theorem~\ref{M-thm:main}}\label{sec:thm1supp}
%======================================================================

\begin{remark}[Several angles near the midpoint]\label{rem:multi}
The bound \eqref{M-eq:dom2} does not require the singular factors to be
treated one cluster at a time: configurations in which several
$\vartheta_j$ approach the midpoint simultaneously are included in the
exact moment $\mathcal M_{N-2}(k+2)$. For a self-contained check of its
finiteness, bound the Vandermonde factor of $p_{N-2}$ by
$4^{(N-2)(N-3)/2}$; then by Tonelli
$\E_{N-2}|\Lambda_{N-2}(1)|^{2a}\le c_N\prod_{j=1}^{N-2}
\int_{\T}|2\sin(\vartheta_j/2)|^{2a}\,d\vartheta_j$, a product of
one-dimensional integrals, each finite iff $2a>-1$. With $2a=4+2k$ this is
$k>-\frac52$.
\end{remark}

\begin{remark}[Uniform integrability]\label{rem:UI}
Lemmas~\ref{M-lem:dom} and~\ref{M-lem:reduce} give, for
$k\in(-\frac52,0]$ and $0<\eps\le1$,
\[
\eps^{-3}\,\E\sum_n\1\{s_n<\eps\}\Big(\frac{M_n}{s_n^2}\Big)^{2k}
\le\frac{5^{-2k}\,\mathcal M_{N-2}(k+2)}{6\pi},
\]
uniformly in $\eps$. Together with Theorem~\ref{M-thm:main}(a) this bounds
$\sup_{\eps\le\eps_0}\langle(M/s^2)^{2k}\rangle_\eps$, so by the de la
Vall\'ee-Poussin criterion $(M_J/s_J^2)^{2k}$ is uniformly integrable
under $\langle\cdot\rangle_\eps$ for every $k>-\frac52$ (use an exponent
$k'\in(-\frac52,k)$). The proof in the main text uses the domination
directly.
\end{remark}

\begin{remark}[Position of the maximum]\label{rem:position}
The maximum of $|\Lambda_N|$ over a small gap sits asymptotically at its
midpoint: if $t_n^*$ is the (smallest) maximiser on the $n$-th gap, then
$\langle\1\{|(t^*-\theta_{(\cdot)})/s-\frac12|>\eta\}\rangle_\eps\to0$ for every
$\eta>0$. Indeed, in the notation of the proof of
Theorem~\ref{M-thm:main}, with $t=sx$, $x\in[0,1]$, the normalised profile
is $x(1-x)\sinc(\frac{sx}2)\sinc(\frac{s(1-x)}2)|\Lambda_{N-2}(e^{isx})|$,
which converges uniformly in $x$ to $|\Lambda_{N-2}(1)|x(1-x)$; when
$\Lambda_{N-2}(1)\ne0$ the limit has the unique maximiser $\frac12$, so the
maximisers converge to $\frac12$. Lemma~\ref{M-lem:reduce} with the bounded
integrand $\1\{|x^*-\frac12|>\eta\}$, the bound $W_s\le4^{2(N-2)}$ and
dominated convergence give the claim.
\end{remark}

%======================================================================
\section{Proof details for Corollary~\ref{M-cor:largeN}}\label{sec:largeNsupp}
%======================================================================

\begin{lemma}[Barnes asymptotics]\label{lem:barnes}
For fixed real $a>-\frac12$, as $n\to\infty$,
\[
\mathcal M_n(a)=\frac{G(1+a)^2}{G(1+2a)}\cdot
\frac{G(n+1)\,G(n+1+2a)}{G(n+1+a)^2}\sim\frac{G(1+a)^2}{G(1+2a)}\,n^{a^2}.
\]
\end{lemma}

\begin{proof}
The identity follows from $\prod_{j=1}^n\Gamma(j+c)=G(n+1+c)/G(1+c)$,
a consequence of $G(z+1)=\Gamma(z)G(z)$. The classical expansion
$\log G(z+1)=\frac{z^2}2\log z-\frac34z^2+\frac z2\log2\pi
-\frac1{12}\log z+\zeta'(-1)+O(z^{-1})$ \cite{Bar00} gives, for fixed $c$,
$\log G(n+1+c)-\log G(n+1)=(cn+\frac{c^2}2)\log n-cn+\frac c2\log2\pi+o(1)$.
The combination with $c=0,2a,a,a$ and weights $1,1,-1,-1$ leaves
$(2a^2-a^2)\log n=a^2\log n$. This is the asymptotic of \cite{KS00}.
\end{proof}

\begin{lemma}[Digamma sums]\label{lem:digamma}
For integers $a\ge0$, $n\ge1$,
\begin{align}
\sum_{j=1}^n\psi'(j+a)&=H_{n+a}-H_a-a\big[\psi'(a+1)-\psi'(n+a+1)\big]
+n\,\psi'(n+a+1),\label{eq:dig1}\\
\sum_{j=1}^\infty\psi^{(r-1)}(j+a)&=(-1)^r(r-1)!\,A_a^{(r)}\qquad(r\ge3).
\label{eq:dig2}
\end{align}
\end{lemma}

\begin{proof}
Write $\psi^{(r-1)}(x)=(-1)^r(r-1)!\sum_{m\ge0}(x+m)^{-r}$ and put
$l=j+a+m$. For fixed $l\ge a+1$ the number of $j\in[1,n]$ with $j\le l-a$
is $\min(n,l-a)$. For $r=2$ this gives
$\sum_{l=a+1}^{a+n}(l-a)l^{-2}+n\sum_{l>a+n}l^{-2}$, which is
\eqref{eq:dig1} since $\sum_{l>c}l^{-2}=\psi'(c+1)$. For $r\ge3$ and
$n=\infty$ it gives $\sum_{l>a}(l-a)l^{-r}=A_a^{(r)}$. All rearrangements
are of series of one sign.
\end{proof}

\emph{Details for (c).} By \eqref{M-eq:cumulants} and \eqref{eq:dig1}
with $n=N-2$, $\kappa_2=S_4-\frac12S_2$ where $S_a=\sum_{j\le n}\psi'(j+a)
=\log n+\gamma-H_a-a(\zeta(2)-H_a^{(2)})+1+O(n^{-1})$, using
$\psi'(a+1)=\zeta(2)-H_a^{(2)}$, $n\psi'(n+a+1)=1+O(n^{-1})$. With
$H_4=\frac{25}{12}$, $H_4^{(2)}=\frac{205}{144}$, $H_2=\frac32$,
$H_2^{(2)}=\frac54$:
$\kappa_2=\frac12\log n+\frac\gamma2+\frac12-\frac{25}{12}+\frac34
-\frac54+\frac{205}{36}-3\zeta(2)+O(n^{-1})
=\frac12\log N+\frac{65}{18}+\frac\gamma2-\frac{\pi^2}2+O(N^{-1})$.

\emph{Details for (d).} For $j=1$,
$\tau_1=-2(\zeta(3)-H_4^{(3)})+\frac12(\zeta(3)-H_2^{(3)})=-0.010261\ldots$.
For $j\ge2$ use $x^{-2}\le|\psi''(x)|\le x^{-2}+2x^{-3}$ (compare
$2\sum_{m\ge0}(x+m)^{-3}$ with $2\int_0^\infty(x+t)^{-3}dt$): it suffices
that $(j+4)^{-2}>\frac14(j+4)(j+2)^{-3}$, i.e.\
$4^{1/3}(j+2)>j+4$, true for $j\ge2$. The limit $\sum_{j\ge1}\tau_j$ is
evaluated by \eqref{eq:dig2}. The general $r$ is the same computation; the
series converge because $|\psi^{(r-1)}(x)|\le(r-1)!x^{-r}+(r-2)!x^{1-r}$.

%======================================================================
\section{How much do small gaps contribute to the second moment?}\label{sec:second}
%======================================================================

\begin{corollary}\label{cor:second}
For $N\ge2$,
\[
\lim_{\eps\to0}\eps^{-7}\,\E\sum_{n=1}^N\1\{s_n<\eps\}M_n^2
=\frac{\mathcal M_{N-2}(3)}{224\pi}
=\frac{(N-1)N^2(N+1)^3(N+2)^2(N+3)}{1\,935\,360\,\pi}.
\]
\end{corollary}

\begin{proof}
$M_n^2=\eps^4(s_n/\eps)^4(M_n/s_n^2)^2$. Apply \eqref{M-eq:main} with the
bounded continuous $F(u,y)=u^4y^2$: the limit is
$\frac1{2\pi}\int_0^1u^6du\cdot\frac1{16}\E_{N-2}|\Lambda_{N-2}(1)|^6
=\mathcal M_{N-2}(3)/(224\pi)$; then use \eqref{M-eq:M23}.
\end{proof}

In unfolded units, $\eps=2\pi\tilde\eps/N$, the contribution of the gaps
with $\tilde s_n<\tilde\eps$ is
\[
\E\sum_{\tilde s_n<\tilde\eps}M_n^2=\Big(\frac{2\pi}N\Big)^7
\frac{\mathcal M_{N-2}(3)}{224\pi}\,\tilde\eps^{\,7}\,(1+o(1)),
\qquad
\Big(\frac{2\pi}N\Big)^7\frac{\mathcal M_{N-2}(3)}{224\pi}\sim
\frac{\pi^6}{15120}\,N^2=0.0636\,N^2 .
\]

\subsection*{The total second moment}
To compare this with the sum over all gaps, let
\[
V_N(\vartheta):=\prod_{j=1}^N2\sin\frac{\vartheta-\theta_j}2 ,
\]
a real, real-analytic function with $|V_N(\vartheta)|=|\Lambda_N(e^{i\vartheta})|$
and $V_N(\vartheta+2\pi)=(-1)^NV_N(\vartheta)$. Winn \cite[(1.1)--(1.2)]{Win12}
uses $Z_U(\vartheta)=\prod_j(1-e^{i(\theta_j-\vartheta)})$ and
$V_U(\vartheta)=\exp\big(iN\frac{\vartheta+\pi}2-\frac i2\sum_j\theta_j\big)
Z_U(\vartheta)$. Since
$1-e^{i(\theta_j-\vartheta)}=2i\,e^{i(\theta_j-\vartheta)/2}
\sin\frac{\vartheta-\theta_j}2$, we have
$Z_U(\vartheta)=i^Ne^{\frac i2(\sum_j\theta_j-N\vartheta)}V_N(\vartheta)$,
hence $V_U=e^{iN\pi/2}\,i^N\,V_N=(-1)^NV_N$. Following Winn, put
$\widetilde F_N(h,k):=\E|V_N(0)|^{2k-2h}|V_N'(0)|^{2h}$.

\begin{proposition}[Total second moment]\label{prop:total}
Let $N\ge2$. Almost surely
\[
\sum_{n=1}^NM_n^2=\int_0^{2\pi}|V_N(\vartheta)V_N'(\vartheta)|\,d\vartheta ,
\]
hence $\E\sum_nM_n^2=2\pi\widetilde F_N(\frac12,1)$, and as $N\to\infty$
\[
\E\sum_{n=1}^NM_n^2\sim\frac{e^2-5}{2}\,N^2=1.194528\ldots N^2 .
\]
\end{proposition}

\begin{proof}
Almost surely the $\theta_j$ are distinct. On the $n$-th gap $V_N$ has no
zero, and $(\log|V_N|)'=\frac12\sum_j\cot\frac{\vartheta-\theta_j}2$ is
strictly decreasing (its derivative is
$-\frac14\sum_j\sin^{-2}\frac{\vartheta-\theta_j}2$), with limits $+\infty$
and $-\infty$ at the left and right endpoints. Hence $V_N^2$ increases up to
a unique point $t_n^*$ and decreases after it, vanishing at both endpoints,
so that $\int_{\mathrm{gap}}|V_NV_N'|=\frac12\int_{\mathrm{gap}}|(V_N^2)'|
=V_N(t_n^*)^2=M_n^2$. The $N$ gaps tile a period of the $2\pi$-periodic
function $|V_NV_N'|$, which gives the identity. Rotating the eigenangles by
$\alpha$ replaces $V_N(\vartheta)$ by $V_N(\vartheta-\alpha)$, so
$\E|V_N(\vartheta)V_N'(\vartheta)|$ does not depend on $\vartheta$ and
$\E\sum_nM_n^2=2\pi\widetilde F_N(\frac12,1)$. Finally, Winn proved
$N^{-2}\widetilde F_N(\frac12,1)\to(e^2-5)/(4\pi)$ \cite[(1.11), proved
in \S7]{Win12}, a conjecture of Hughes.
\end{proof}

The identity is the random-matrix counterpart of the reduction used by
Conrey and Ghosh for Hardy's function (see \cite[\S1]{Win12}), and the
constants match: under the Riemann Hypothesis
$\sum_{0<\gamma\le T}\max_{\gamma\le t\le\gamma^+}|\zeta(\frac12+it)|^2\sim
\frac{e^2-5}{4\pi}T(\log T)^2$ \cite{CG85}, i.e.\ a mean of
$\frac{e^2-5}2\log T$ per gap, against $\frac{e^2-5}2N$ per gap here.
Winn's Theorem~2.2 \cite{Win12} also evaluates $\widetilde F_N(\frac12,1)$
at finite $N$; for $k=1$, $h=\frac12$ it reduces to (see below)
\begin{equation}\label{eq:winnN}
\E\sum_{n=1}^NM_n^2=2(N+1)\Big[N-\sum_{p=2}^N\frac{(p-2)!\,2^p\,N!}
{(N-p)!\,(p+1)!\,p!\,N^{p-1}}\Big]
\end{equation}
(for $N=2$ this is $10$, as a direct computation confirms), so that
$\E\sum_nM_n^2/N^2=1.6055$, $1.4381$, $1.3457$, $1.2949$, $1.2696$ for
$N=6,10,16,24,32$. A Monte Carlo run ($4\times10^4$ matrices for $N\le16$,
$2\times10^4$ for $N=24,32$; maxima on the $63$ interior points of a
$65$-point grid per gap with parabolic refinement) gives $1.604(5)$,
$1.455(7)$, $1.339(8)$, $1.284(14)$, $1.258(14)$; the $N=10$ value is
$2.6\sigma$ high, and an independent run with $4\times10^5$ matrices gives
$1.4394(20)$ there.

Gaps below $\tilde\eps$ therefore carry a fraction
$c_N\tilde\eps^{\,7}(1+o(1))$ of the second moment, with $c_N=0.140$
($N=6$) decreasing to $0.066$ ($N=32$) and
$c_N\to\pi^6/(7560(e^2-5))=0.0532$; this is about $10^{-6}$ for
$\tilde\eps=0.2$. Small gaps are invisible in the mean of $M_n^2$, which is
the quantity controlled in \cite{CG85,HLPC24}.

\subsection*{Winn's formula for $h=\frac12$, $k=1$}\label{sec:winn}
We specialise \cite[Thm.~2.2, eq.~(2.9)]{Win12} to $h=\frac12$, $k=1$, in
Winn's notation. For a partition $\lambda$ and $b\in\mathbb R$,
$[b]_\lambda=\prod_{i=1}^{\ell(\lambda)}(b-i+1)_{\lambda_i}$ with the rising
factorial $(x)_m=x(x+1)\cdots(x+m-1)$, $h_\lambda$ is the product of the
hook lengths of $\lambda$, and
\[
C_N(p,k)=(-2)^p\sum_{\lambda\vdash_kp}\frac{[k]_\lambda\,[-N]_\lambda}
{[2k]_\lambda\,h_\lambda^2},
\]
the sum running over partitions of $p$ into at most $k$ parts
\cite[(2.1)--(2.5)]{Win12}. For half-integer $h$ and $k\in\mathbb N$,
\cite[(2.9)]{Win12} reads
\begin{multline*}
\widetilde F_N(h,k)=\frac{2(-1)^{h+1/2}}{2^{2h}\pi}\,\widetilde F_N(0,k)
\Bigg\{\sum_{p=1}^{2h}\sum_{\ell=1}^p\binom{2h}{p-\ell}\frac{(-1)^\ell}{\ell}
(-N)^{2h-p}\,p!\,C_N(p,k)\\
+\sum_{p=2h+1}^{kN}\frac{(2h)!\,(p-2h-1)!}{N^{p-2h}}\,C_N(p,k)\Bigg\}.
\end{multline*}

\emph{Step 1: the ingredients.} For $h=\frac12$, $k=1$: the prefactor is
$2(-1)^1/(2\pi)=-1/\pi$; $\widetilde F_N(0,1)=\mathcal M_N(1)=\prod_{j=1}^N
\frac{\Gamma(j)\Gamma(j+2)}{\Gamma(j+1)^2}=\prod_{j=1}^N\frac{j+1}j=N+1$ by
\eqref{M-eq:KS}; the first double sum has the single term $p=\ell=1$, equal
to $\binom10\frac{(-1)^1}1(-N)^0\,1!\,C_N(1,1)=-C_N(1,1)$; and the second
sum runs over $2\le p\le N$ with weight $(p-2)!/N^{p-1}$.

\emph{Step 2: the partition sum.} For $k=1$ only the one-part partition
$\lambda=(p)$ contributes. Its hook lengths are $p,p-1,\dots,1$, so
$h_{(p)}=p!$, and $[b]_{(p)}=(b)_p$, so that $[1]_{(p)}=p!$,
$[2]_{(p)}=(p+1)!$ and $[-N]_{(p)}=(-N)(-N+1)\cdots(-N+p-1)=(-1)^pN!/(N-p)!$
for $p\le N$. Hence
\[
C_N(p,1)=(-2)^p\,\frac{p!\,(-1)^pN!/(N-p)!}{(p+1)!\,(p!)^2}
=\frac{2^p\,N!}{(N-p)!\,(p+1)!\,p!},\qquad C_N(1,1)=N .
\]

\emph{Step 3: assembly.} Collecting,
\[
\widetilde F_N\big(\tfrac12,1\big)=\frac{N+1}\pi\Big[N-\sum_{p=2}^N
\frac{(p-2)!\,2^p\,N!}{(N-p)!\,(p+1)!\,p!\,N^{p-1}}\Big],
\]
and $\E\sum_nM_n^2=2\pi\widetilde F_N(\frac12,1)$ (Proposition~\ref{prop:total})
gives \eqref{eq:winnN}. Each term is rational, so $\E\sum_nM_n^2\in\mathbb Q$;
for $N=2,3,4,6$ it equals $10$, $\frac{496}{27}$, $\frac{233}8$,
$\frac{42134}{729}$.

\emph{Checks.} (i) For $N=2$ the bracket is $2-\frac13$, so the formula
gives $10$. Directly: let $\phi\in(0,2\pi)$ be $\theta_2-\theta_1$ modulo
$2\pi$; the two gaps are $\phi$ and $2\pi-\phi$, and by
the proof of Lemma~\ref{M-lem:hump} (with $N-2=0$ other angles) a gap of length $s$ has
hump $4\sin^2\frac s4$, attained at its midpoint. The humps are thus
$4\sin^2\frac\phi4$ and $4\cos^2\frac\phi4$, so
$\sum_nM_n^2=16(\sin^4\frac\phi4+\cos^4\frac\phi4)=16-8\sin^2\frac\phi2$,
and $\phi$ has density $\sin^2(\phi/2)/\pi$ on $(0,2\pi)$ by
\eqref{M-eq:weyl}, so $\E\sin^2\frac\phi2=\frac34$ and $\E\sum_nM_n^2=10$.
(ii) As $N\to\infty$, $N!/((N-p)!N^p)\to1$ for each $p$ and the terms are
dominated by $2^p(p-2)!/((p+1)!p!)$, so
$N^{-2}\widetilde F_N(\frac12,1)\to\frac1\pi\big[1-\sum_{p\ge2}
\frac{2^p(p-2)!}{p!(p+1)!}\big]=\frac{e^2-5}{4\pi}$, which is
\cite[(7.8)]{Win12}.

%======================================================================
\section{Monte Carlo verification in full}\label{sec:mcsupp}
%======================================================================

\subsection*{Method}
Haar unitaries were generated by the QR decomposition of complex Ginibre
matrices with the phases of the diagonal of $R$ removed \cite{Mez07}. For
every gap with unfolded length $\tilde s<0.4$, $\log|\Lambda_N|$ was
evaluated on the $39$ interior points of a uniform $41$-point grid across
the gap and the maximum refined by a parabola through the largest grid
value and its neighbours. Statistics are reported in the bands
$\tilde s<\tilde\eps$ for $\tilde\eps\in\{0.4,0.3,0.2,0.1,0.05\}$, in
angle units ($M/s^2$); the predictions are Theorem~\ref{M-thm:main} with
Theorem~\ref{M-thm:product}. The ``rate'' is
$\#\{\text{gaps}<\eps\}/(\text{matrices}\cdot\eps^3)$, to be compared with
$N^2(N^2-1)/(72\pi)$. No zeta data enter this section.

\begin{table}[h]
\centering
\small
\begin{tabular}{rrrcccccc}
\toprule
& & & \multicolumn{3}{c}{$\E[M/s^2]$} & \multicolumn{3}{c}{rate}\\
\cmidrule(lr){4-6}\cmidrule(lr){7-9}
$N$ & matrices & events & Thm & $\tilde s<0.2$ & $\tilde s<0.1$
& Thm & $\tilde s<0.2$ & $\tilde s<0.1$\\
\midrule
6 & $6\cdot10^5$ & 3\,724 & 1.4784 & 1.4828(30) & 1.4751(83) & 5.570 & 5.392 & 5.405\\
10 & $4\cdot10^5$ & 4\,302 & 3.9279 & 3.935(11) & 3.919(31) & 43.77 & 42.96 & 43.36\\
16 & $3\cdot10^5$ & 5\,253 & 10.194 & 10.178(35) & 10.059(98) & 288.6 & 280.4 & 289.1\\
24 & $2\cdot10^5$ & 5\,229 & 23.903 & 24.13(10) & 24.10(29) & 1464 & 1419 & 1457\\
\bottomrule
\end{tabular}

\medskip
\begin{tabular}{rccccccc}
\toprule
& \multicolumn{3}{c}{$\E\log(M/s^2)$} & \multicolumn{3}{c}{$\Var\log(M/s^2)$} & $\kappa_3$\\
\cmidrule(lr){2-4}\cmidrule(lr){5-7}\cmidrule(lr){8-8}
$N$ & Thm & $\tilde s<0.2$ & $\tilde s<0.1$ & Thm & $\tilde s<0.2$ & $\tilde s<0.1$ & Thm\\
\midrule
6 & 0.3232 & 0.3276(22) & 0.3229(62) & 0.1486 & 0.1447 & 0.1448 & $-0.0439$\\
10 & 1.2292 & 1.2323(29) & 1.2330(81) & 0.2987 & 0.2963 & 0.2834 & $-0.0688$\\
16 & 2.1007 & 2.0998(34) & 2.0822(95) & 0.4691 & 0.4664 & 0.4781 & $-0.0872$\\
24 & 2.8723 & 2.8827(39) & 2.8804(110) & 0.6340 & 0.6306 & 0.6326 & $-0.0990$\\
\bottomrule
\end{tabular}
\caption{\label{tab:mcfull}CUE Monte Carlo against
Theorems~\ref{M-thm:main} and~\ref{M-thm:product} (angle units;
statistical errors in the last digits; ``events'' is the number of gaps
with $\tilde s<0.1$). The lower half is Table~\ref{M-tab:mc}. Files
\texttt{200t\_mc\_N\{6,10,16,24\}.json}.}
\end{table}

\subsection*{Results}
Table~\ref{tab:mcfull} shows agreement with the theorems at the level of the
statistical errors: at $\tilde s<0.1$ every deviation of the mean hump and
of the log-mean is below $2\sigma$ and the counts agree within $2\sigma$
(Poisson errors). At $\tilde s<0.2$ the small finite-$\eps$ bias
discussed below is visible (largest in the log-mean at $N=24$, $+0.010$).
Further checks:
\begin{itemize}
\item \emph{Independence.} The correlation of $\tilde s$ with
$\log(M/s^2)$ falls with the band: $0.021$, $0.013$, $0.007$, $0.000$ at
$\tilde\eps=0.4,0.3,0.2,0.1$ ($N=10$, file \texttt{200t\_mc\_N10.json}).
\item \emph{Lemma~\ref{M-lem:hump}.} The largest difference between
$\log(M/s^2)$ and $\log(|\Lambda_{N-2}(\mathrm{mid})|/4)$ over all events
with $\tilde s<0.1$ is at most $0.042$ (all $N$; files of
Table~\ref{tab:mcfull}).
\item \emph{Third cumulant.} Independent runs with $8\times10^5$ ($N=10$)
and $4\times10^5$ ($N=16$) matrices give $\kappa_3=-0.0639$ ($8\,526$
events) and $-0.0863$ ($7\,021$ events) at $\tilde s<0.1$, against
$-0.0688$ and $-0.0872$ (files \texttt{200t\_mc3\_N\{10,16\}.json}); the same
runs give $\E\log(M/s^2)=1.2357(59)$ and $2.1021(81)$.
\item \emph{Theorem~\ref{M-thm:product} directly.} Sampling the $\xi_j$ by
rejection ($4\times10^5$ products) gives, for $N=10$, $\E\log(X/4)=1.2286$
($1.2292$), $\Var=0.2976$ ($0.2987$), $\kappa_3=-0.0675$ ($-0.0688$),
$\E X/4=3.924$ ($3.928$); at $N=24$ the mean, variance and $\E X/4$ agree
within $0.2\%$, and $\kappa_3=-0.0963$ ($-0.0990$) (script
\texttt{200t\_carpim.py}, output \nolinkurl{not8_ek/200t_carpim_seed7.out}).
\item \emph{Finite $\eps$.} In a run with $3\times10^6$ matrices at $N=10$
(file \nolinkurl{not8_ek/200t_mc_N10_3e6_seed7.json})
the excess of $\E[M/s^2]$ over its limit is $0.0771(15)$, $0.0440(23)$,
$0.0198(41)$ at $\tilde\eps=0.4$, $0.3$, $0.2$: the successive ratios $1.75$
and $2.22$ match $(4/3)^2=1.78$ and $(3/2)^2=2.25$, as expected from
Remark~\ref{M-rem:finiteeps}. At $\tilde\eps=0.2$ the excesses of
$\E\log$, $\Var\log$ and $\kappa_3$ are $+0.0070(11)$, $-0.005$ and
$+0.004$.
\end{itemize}

%======================================================================
\section{The zeta side before the tests}\label{sec:zetasupp}
%======================================================================

\subsection*{Weak points of Conjecture~\ref{M-conj:zeta}}
Recorded before the tests: (W1) the independence of $P_X$ from the
conditioning, whereas the prime waves modulate the local density of zeros
at relative order $1/L$ (Landau's formula \cite{Landau12,Gonek93}) and
imprint on the gaps, as measured in the companion notes
\cite{companion2,companion3} --- and at $b=2$ the arithmetic term is not a
small correction ($\kappa_3$: CUE baseline $-0.069$, arithmetic shift
$+0.234$ at $L\approx10$); (W2) finite $L$, where \cite{companion1} found
effective dimensions $N=L+c$; (W3) the order of limits and the transfer of
finite-$\eps$ corrections; (W4) the passage from (ii) to (iii), which
interchanges $L\to\infty$ with differentiation in $k$ at $k=0$.

\subsection*{Finite $L$}
At $L\approx11$ the form $a_k\times$(tilted CUE) is not the moment sequence
of a probability law: $\log|a_{i\omega/2}|$ grows faster than $\omega^2$
(each prime $p\lesssim\omega^2$ contributes), while the characteristic
function of the tilted CUE decays only like
$e^{-\omega^2\log N/4}$. Its low cumulants are nevertheless defined.
Genuine probability models are the hybrids $\mathrm{hyb}(X)$: an
independent random Euler product over $p\le X$ times a tilted CUE of
dimension $N_X=L/(e^\gamma\log X)$. Table~\ref{tab:models} lists the
predictions at $L=10$. At $b=2$ every arithmetic model predicts a positive
$\kappa_3$, whereas the pure tilted CUE gives $-0.069$ and a stationary
Gaussian surface with the same power spectrum gives
$\kappa_3=\psi''(3/2)/8=-0.1036$ and $\kappa_2=\psi'(3/2)/4=0.2337$,
independent of $L$ (the two-point Kac--Rice density makes
$|X''|/\sigma$ a $\chi_3$ variable). The pre-test headline was therefore a
sign: \emph{close-pair log-humps of zeta are skewed to the right}, with
skewness $\kappa_3/\kappa_2^{3/2}\approx+1.7$ at $L=10$ under $a_k$.

\begin{table}[h]
\centering
\small
\begin{tabular}{ccccccc}
\toprule
$b$ & CUE & $a_k$ & hyb $X{=}2$ & $X{=}3$ & $X{=}5$ & $X{=}7$\\
\midrule
0 & $1.940/{-2.393}$ & $1.852/{-2.159}$ & $2.126/{-2.233}$ & $2.080/{-2.126}$ & $1.996/{-2.042}$ & $1.977/{-1.991}$\\
1 & $0.568/{-0.194}$ & $0.480/{+0.039}$ & $0.775/{-0.039}$ & $0.790/{+0.050}$ & $0.778/{+0.103}$ & $0.803/{+0.131}$\\
2 & $0.299/{-0.069}$ & $0.211/{+0.165}$ & $0.523/{+0.083}$ & $0.585/{+0.160}$ & $0.622/{+0.197}$ & $0.674/{+0.213}$\\
\bottomrule
\end{tabular}
\caption{\label{tab:models}Pre-test predictions $\kappa_2/\kappa_3$ of the
log-observable on rung $b$ at $L=10$. Gaussian surface: $b=0$:
$\pi^2/8=1.234$, $-7\zeta(3)/4=-2.104$; $b=1$: $\pi^2/24=0.411$,
$-\zeta(3)/4=-0.300$; $b=2$: $0.234$, $-0.104$.}
\end{table}

%======================================================================
\section{The pre-registered tests in full}\label{sec:testsupp}
%======================================================================

\subsection*{Protocol}
Each test was frozen --- predictions, thresholds, analysis code --- in a
pre-registration file that was committed with its sha256 stamp to the
author's private working repository, and pushed, before the measurement:
200-A on 26 September 2026, 17:12 (UTC$+3$), 200-B on 26
September 2026, 20:19, 200-C on 27 September 2026, 16:22. Before each
measurement no distribution summary (moment, cumulant, quantile,
histogram) of the tested observable had been computed on the data used;
engine checks reported only differences between two engines. (The only
earlier look at a rung-$0$ distribution in the project was a
Kolmogorov--Smirnov test of $\log|Z(t)|$ against Selberg's normal law at
low height; $Z'$ at the zeros and the humps of close pairs had never been
computed.) Test 200-B was designed after 200-A was sealed, and 200-C after
200-B. Test 201 (phase lock) was frozen in the same way on 28 September
2026, 09:34, and committed and pushed to the (private) repository before
its measurement; no phase statistic of the zeros had been computed before
(only synthetic data).

\subsection*{Data and engine}
Tests 200-A/B use A.~M.~Odlyzko's first $2\times10^6$ zeros
\cite{OdlyzkoTables,Odl87} in three windows: W1, zeros
$20\,868$--$198\,239$ ($L\in[8,10)$, $\bar L\approx9.3$); W2,
$198\,239$--$598\,742$ ($[10,11)$, $10.6$); W3,
$598\,742$--$2\,001\,052$ ($[11,12.1]$, $11.65$). Test 200-C uses four
LMFDB files of zeros computed by Platt \cite{LMFDB,Pla15}, $1.5\times10^6$
consecutive zeros each, at $\bar L=14.19$, $16.58$, $18.88$, $22.31$
(windows C1--C4). $Z$ and $Z'$ are computed by a vectorised
Riemann--Siegel formula with the $C_0$ and $C_1$ remainder terms (anchored
phases in C1--C4), checked against \texttt{mpmath} ($|\Delta Z|\le
2.4\times10^{-7}$ in W1--W3, $\sim10^{-9}$ in C1--C4). On rung $b=0$ we use
$10^6$ uniform random $t$ per window; on $b=1$ all zeros (in C4, $10^6$
random zeros); on $b=2$ all pairs with $\tilde\delta_n<\eps$, primary
$\eps=0.2$ ($1\,143$, $2\,860$, $10\,165$ pairs in W1--W3; $11\,515$ to
$12\,143$ in each of C1--C4), with $M_n$ from a $129$-point grid and
parabolic refinement. Statistics are unbiased $k$-statistics $k_2,k_3$;
errors are jackknife errors over $64$ contiguous $t$-blocks (128 where a
pre-registered block-sensitivity rule required it), rescaled by factors
$f$ calibrated on synthetic data processed by the same pipeline
($f=0.82$--$1.22$ in 200-A, $0.94$--$1.13$ in 200-B, $0.94$--$1.25$ in
200-C).

\subsection*{Shifts}
The \emph{arithmetic shift} of rung $b$ is
$s_b:=k_3-\kappa_3^{\mathrm{CUE},b}(\bar L)$ (similarly for $k_2$). On
$b=2$ we also subtract the finite-$\eps$ correction
$\Delta\kappa_r^{\mathrm{CUE}}(\eps)$ from a CUE Monte Carlo at the same
$\eps$ and $N\approx\bar L$, and add $100\%$ of it to the error, since its
transfer to zeta is assumed, not tested; this defines
$D_r:=k_r-\kappa_r^{\mathrm{CUE},2}(\bar L)-\Delta\kappa_r^{\mathrm{CUE}}(\eps)$,
pooled over windows by inverse variance ($\bar D_r$).

\subsection*{Test 200-A: the lower rungs}
A family of parameter-free models --- CUE, $a_k$, $\mathrm{hyb}(X)$ for
$X\in\{2,3,5,7,11\}$, Gaussian --- was compared on the twelve numbers
($k_2,k_3$ on two rungs in three windows) with the full jackknife
covariance plus a pre-declared theory margin ($0.0188$ for $\kappa_2$,
$0.0334$ for $\kappa_3$).

\begin{observation}[Lower rungs; pre-registered]\label{obs:A}
The arithmetic signature is present ($\chi^2_{\mathrm{CUE}}-\chi^2_{a_k}=287.1$,
threshold $25$), and the selected model is $a_k$, decisively (runner-up
CUE at $\Delta\chi^2=287.1$, threshold $9$), with goodness of fit
$\chi^2_{\min}/12=28.2/12=2.35$ (``good''). The two rungs separately select
$a_k$ ($\chi^2=6.1$ and $22.1$). The hybrids and the Gaussian surface are
rejected ($\chi^2$: CUE $315.2$, $a_k$ $28.2$, hybrids $872$--$1243$,
Gaussian $3689.7$; a continuous-$X$ hybrid is best at $X=2.99$ with
$\chi^2=129.3$).
\end{observation}

Post hoc, the shifts tell a more detailed story (Table~\ref{tab:ladder}).
The variance shift is nearly rung-independent and close to
$\kappa_2^{\mathrm{arit}}$: $-0.0820\pm0.0024$ at $b=0$ and
$-0.0726\pm0.0008$ at $b=1$, against $-0.0881$. The skewness shift is
not: $+0.283\pm0.009$ at $b=0$ and $+0.150\pm0.001$ at $b=1$, against
$+0.2337$, stable across the three windows. A ladder with the same
arithmetic contribution on every rung therefore fails at the level of
$\kappa_3$ at these heights. Before 200-B we registered, besides $a_k$,
two extrapolations of the two lower rungs (not theories): geometric
decay, $s_2=s_1^2/s_0$, predicting $\bar D_3=+0.0794\pm0.0028$; and linear
decay, $s_2=2s_1-s_0$, predicting $+0.017\pm0.0095$ (grouped with pure
RMT, $\bar D_3=0$, as ``extinguished'').

\subsection*{Test 200-B: close pairs}

\begin{observation}[Close pairs; pre-registered, blind]\label{obs:B}
With $14\,168$ close pairs ($\tilde\delta<0.2$, $L\approx9$--$12$):
\begin{enumerate}
\item[(1)] The arithmetic signature is present on $b=2$:
$\bar D_3=+0.0859\pm0.0030$ ($29\sigma$).
\item[(2)] The sign of the skewness flips: the pooled corrected
$\kappa_3=+0.0143\pm0.0030$ ($4.8\sigma$). The skewness
$\kappa_3/\kappa_2^{3/2}$ is about $+0.15$, not the $+1.7$ of the $a_k$
form: the headline prediction held in sign, not in size.
\item[(3)] The $a_k$ model carried over from 200-A holds for the variance
($\bar D_2=-0.0988$ against $-0.0881$, tolerance $0.024$) and fails for
the third cumulant (difference $-0.148$, tolerance $0.142$): partial.
\item[(4)] The shape of the ladder is geometric, decisively:
$\chi^2=0.36$ (geometric), $24.0$ (extinguished), $200.7$ (constant $a_k$).
The blind geometric prediction $+0.0794\pm0.0028$ differs from the
measurement by $0.0065$ ($1.6\sigma$).
\end{enumerate}
\end{observation}

Secondary results (recorded, not decisive): $\bar D_3=+0.0867\pm0.0066$ at
$\eps=0.1$ and $+0.0828\pm0.0044$ at $\eps=0.3$; per window $+0.076(6)$,
$+0.084(5)$, $+0.093(5)$, a mild rise with height ($\approx2.5\sigma$). The
variance shift on $b=2$ is $-0.0988\pm0.0036$, far from the geometric and
linear extrapolations ($-0.064$, $\approx9\sigma$).
Theorem~\ref{M-thm:main}(b) has zeta counterparts: the Pearson correlation
of $\tilde\delta_n$ and $\log Y_n$ is $+0.021$, $-0.004$, $-0.005$ in
W1--W3 (CUE Monte Carlo: $+0.005$ to $+0.008$), and the fraction of events
with $\tilde\delta_n<\eps/2$ is $0.124$--$0.129$ (the $U^{1/3}$ law gives
$0.125$, the CUE at finite $\eps$ $0.128$).

\begin{table}[h]
\centering
\small
\begin{tabular}{llccc}
\toprule
rung & observable & $\kappa_3$ shift & ratio to previous rung & $\kappa_2$ shift\\
\midrule
$b=0$ & $\log|Z(t)|$, random $t$ & $+0.2828\pm0.0093$ & --- & $-0.0820\pm0.0024$\\
$b=1$ & $\log|Z'(\gamma_n)|$ & $+0.1499\pm0.0009$ & $0.530\pm0.018$ & $-0.0726\pm0.0008$\\
$b=2$ & $\log(M_n/\tilde\delta_n^2)$, $\tilde\delta_n<0.2$ & $+0.0859\pm0.0030$ & $0.573\pm0.020$ & $-0.0988\pm0.0036$\\
\midrule
$a_k$ & (every rung) & $+0.2337$ & $1$ & $-0.0881$\\
\bottomrule
\end{tabular}
\caption{\label{tab:ladder}Arithmetic shifts (observed minus tilted CUE at
$N=\bar L$), pooled over W1--W3 ($L\approx9$--$12$).}
\end{table}

\subsection*{Test 200-C: height}
Is the decay a property of the ladder or of the height? Before 200-C, a
post hoc fit of the A/B window values to $s_b(L)=a_k+c_bL^{-\nu}$ gave
$\nu=0.48\pm0.12$; three hypotheses were registered for the new windows:
$H_C$ (shifts constant at their A/B values), $H_{S\nu}$ (the fitted
relaxation towards $a_k$) and $H_{S1}$ ($\nu=1$). The exponent is called
$\gamma$ in the pre-registration files.

\begin{observation}[Height; pre-registered, blind]\label{obs:C}
On C1--C4 ($L=14.2$--$22.3$), the class is $H_{S\nu}$, decisively
($\chi^2=10.3$ for $12$ points; $H_{S1}$ $30.1$, $H_C$ $166.2$). Fitted on
the new windows alone with asymptote $a_k$, the exponent is
$\hat\nu=0.749\pm0.170$ ($4.4\sigma$ from zero: finite height). Over all
seven windows with a free common asymptote, $\hat A=0.277\pm0.041$
(within $1.0\sigma$ of $\kappa_3^{\mathrm{arit}}=0.2337$) and
$\hat\nu=0.447\pm0.144$.
\end{observation}

The per-window shifts are in Table~\ref{M-tab:shifts}. The blind
$H_{S\nu}$ predictions for C1--C4 were: $b=1$: $+0.159$, $+0.164$,
$+0.169$, $+0.174$; $b=2$: $+0.105$, $+0.114$, $+0.121$, $+0.130$. The
observed $b=1$ shifts lie $0.005$--$0.01$ above them (faster relaxation).
The variance shifts do not depend on height: $-0.082$ to $-0.094$ ($b=0$),
$-0.075$ to $-0.079$ ($b=1$), $-0.095$ to $-0.115$ ($b=2$), around
$\kappa_2^{\mathrm{arit}}=-0.088$. The gap--hump correlation agrees with
the CUE Monte Carlo ($|\rho|\le0.01$).

\subsection*{Test 201: phase lock at the conditioned points}
The ladder concerns the size of $|Z|$ at conditioned points; a simpler
question is where the conditioned points sit relative to the phase of a
prime wave $\cos(t\log p)$. For single zeros the answer is Landau's formula
\cite{Landau12} (uniform version in \cite{Gonek93}): for fixed $x>1$,
$\sum_{0<\gamma\le T}x^{\rho}=-\frac T{2\pi}\Lambda(x)+O(\log T)$, which for
zeros on the critical line and a window at height $t$ means
\[
\E\big[\cos(m\gamma_n\log p)\big]\approx-\frac{\log p}{p^{m/2}L},\qquad
\E\big[\sin(m\gamma_n\log p)\big]\approx0 .
\]
Heuristically, the explicit formula modulates the local density of zeros,
$n(t)\approx\frac L{2\pi}\big[1-\frac2L\sum_{p,m}\frac{\log p}{p^{m/2}}
\cos(mt\log p)\big]$, and the zeros avoid the crests of the prime waves;
averaging $\cos(mt\log p)$ against $n(t)$ gives the same values. For the
midpoints of close pairs we registered three classes for the ratio $R_2$
of their lock to the single-zero value: $1$ (a pair behaves like one
zero), $2$ (the two densities multiply) and $4$. The last is a heuristic,
not a theorem: if the zeros are locally GUE at density $n(t)$, the density
of pairs with gap $\delta$ at midpoint $t$ is
$n(t)^2[1-\sinc^2(\pi\delta n(t))]\approx\frac{\pi^2}3\delta^2n(t)^4$, and
since the pairs are selected by $\tilde\delta_n=\delta_nL/2\pi<\eps$, i.e.\
with the mean rather than the local density, their midpoints are weighted
by $n(t)^4$ and lock four times as strongly as single zeros; for larger
$\eps$ the weight tends to $n(t)^2$, and $R_2$ towards $2$.

On the seven windows W1--W3, C1--C4 and the primes $p\le13$ (ratios
pooled over primes with the jackknife covariance, then over windows):
\begin{observation}[Phase lock; pre-registered]\label{obs:phase}
\begin{enumerate}
\item[(1)] Single zeros follow Landau's formula to about four decimals:
the ratio of the measured lock to $-\log p/(\sqrt p\,L)$ is
$R_1=1.0000\pm0.00004$ (per window from $1.003\pm0.007$ in W1 to
$1.00003\pm0.00009$ in C4), also at the second harmonic
($1.0000\pm0.0001$), with no drift in $\log L$; at C4, for $p=2$, the mean
of $\cos(\gamma_n\log2)$ is $-0.021975$ against $-0.021973$. At random $t$
the same statistic is $-0.010$ to $+0.005$ (control).
\item[(2)] Midpoints of close pairs ($\tilde\delta_n<0.2$) lock
$R_2=3.874\pm0.023$ times as strongly: class $4$, with the classes $1$ and
$2$ excluded. The ratio decreases with $\eps$, $R_2=4.07\pm0.06$,
$3.87\pm0.02$, $3.74\pm0.01$ at $\eps=0.1,0.2,0.3$, as the heuristic
expects; at the second harmonic it is $3.06\pm0.05$, consistent with the
second-order terms of the $n(t)^4$ weight (a synthetic process with that
weight gives $2.90$).
\end{enumerate}
\end{observation}

For single zeros this verifies a theorem rather than discovers anything; a
by-product is that the block-to-block fluctuations of these averages are
$7$--$19$ times smaller than for independent points, a sign of the
rigidity of the zeros. The lock of close pairs has the size the local-GUE
heuristic predicts, so there is no excess lock. It therefore does not by
itself explain the decay of the $\kappa_3$ shift along the ladder: a
translation of the measured lock into cumulant shifts through independent
Euler factors was inconclusive (its sign depends on the harmonic
cut-off), and the mechanism of the decay remains open. A phase bias of the
primes at the conditioned points, which weakens like $1/L$, would produce
a relaxation of the kind seen in 200-C; as in \cite{companion7}, where a
post hoc pattern $\eta_p\approx p^{-1/3}$ of transfer factors turned out to
be a finite-height effect, a pattern seen at one height ($\approx0.55$ per
rung) was a coincidence of height.

%======================================================================
\section{Numerical details of the ratios-conjecture computations}\label{sec:numsupp}
%======================================================================

\emph{Rung $b=0$, $\kappa_3$.} We use primes up to $10^6$ with integral
tails, Gauss--Legendre panels adapted to the scales $x_1$ and $1/L$
(refining from $16$ to $24$ nodes changes $\kappa_3$ by $6\cdot10^{-4}$),
and drop the second swap term for $w>0.75$, where it carries $e^{-0.75L}$
(changing the cut to $0.5$ moves $\kappa_3$ by $6\cdot10^{-5}$).

\emph{Rung $b=1$, variance.} Near $a=0$ terms of order $1/a$ cancel; we
integrate from $a=0.01/L$, adding the remaining strip linearly (the inner
integral vanishes linearly there). For $\zeta$, $12$, $16$ and $20$ nodes
per panel agree to $5\cdot10^{-6}$, and the swap cut changes the result by
at most $6\cdot10^{-4}$ (at $L=9.35$).

\emph{Rung $b=1$, third cumulant.} In $\kappa_3=\kappa_3^{\rm
CUE}(L)+\iiint(I_\zeta-I_{\rm CUE})$ the difference of the integrands is
integrated over shifts $\ge0.0075$, with the thin strips below added
linearly (cuts from $0.005$ to $0.01$ change $\kappa_3$ by
$5\cdot10^{-5}$, the grid by $10^{-4}$). The interpolation of $A_{\rm
dbl}$ outside its region of convergence moves $\kappa_3$ by at most
$3\cdot10^{-4}$.

%======================================================================
\section{Reproducibility: file guide}\label{sec:repro}
%======================================================================

All code, data products, pre-registration files (with sha256 stamps and
dates) and the research log, including the failed predictions, are
available in the repository
\url{https://github.com/azizran/riemann-zero-statistics} (directory
\texttt{qm\_riemann/}), archived at Zenodo,
\href{https://doi.org/10.5281/zenodo.22942475}{doi:10.5281/zenodo.22942475}.
The commit hashes below refer to the author's private working repository,
in which the pre-registrations and predictions were committed before the
measurements; its history can be shown to referees on request. Paths below are relative to \nolinkurl{qm_riemann/200_configs/}, whose
subdirectory \texttt{not8\_ek/} holds the runs added for this note.
\begin{itemize}
\item Tables~\ref{M-tab:mc} and~\ref{tab:mcfull}, independence and
Lemma~\ref{M-lem:hump} checks: \nolinkurl{200t_mc_N*.json} ($N=6,10,16,24$;
script \nolinkurl{200t_mc.py}, predictions \nolinkurl{200t_tahmin.py});
third cumulant: \nolinkurl{200t_mc3_N*.json} ($N=10,16$); finite $\eps$:
\nolinkurl{not8_ek/200t_mc_N10_3e6_seed7.json}
(arguments \texttt{10 3000000 7}); Theorem~\ref{M-thm:product} by
sampling and $\kappa_2^{\mathrm{arit}}$:
\nolinkurl{200t_carpim.py}, output
\nolinkurl{not8_ek/200t_carpim_seed7.out}.
\item Proposition~\ref{prop:total}: Monte Carlo
\nolinkurl{not8_ek/m2total.py}, output
\nolinkurl{not8_ek/m2total_seed11.out}; independent run
\nolinkurl{not8_ek/m2total_N10_4e5_seed2026.json}; exact values
from Winn's formula \nolinkurl{not8_ek/m2_exact_winn.json} (these
two JSON files contain their generating scripts).
\item Constants, ladder baselines and Table~\ref{tab:models}:
\nolinkurl{200t_sabitler.py}, \nolinkurl{200t_merdiven.py},
\nolinkurl{200t_hibrit.py}, \nolinkurl{200t_aritmetik_kesin.py}.
\item Tests 200-A/B/C: pre-registrations
\nolinkurl{ONKAYIT_200A/B/C.json} (sha256 prefixes
\texttt{f0f0ce57}, \texttt{b37dc244}, \texttt{391257a4}; commits
\texttt{aa861b5}, \texttt{85f0489}, \texttt{567af03}), verdicts
\nolinkurl{HUKUM_200A/B/C.json} (which also hold the CUE Monte Carlo
comparisons of Section~\ref{sec:testsupp}), synthetic calibrations and
finite-$\eps$ CUE corrections \nolinkurl{M6_200A.json},
\nolinkurl{M6b_200B.json}, \nolinkurl{M6c_200C.json}. Test 201: directory
\nolinkurl{qm_riemann/201_configs/}, pre-registration
\nolinkurl{ONKAYIT_201.json} (sha256 prefix \texttt{ecc7cf75}, commit
\texttt{16e7c3b}), verdict \nolinkurl{HUKUM_201.json}.
\item Table~\ref{M-tab:ratios}, Section~\ref{M-sec:ratios} and
Appendix~\ref{M-app:ratios}: directory
\nolinkurl{qm_riemann/203_configs/} (\nolinkurl{203_oran_k2.py},
\nolinkurl{203_oran_k3.py}, \nolinkurl{203_oran_b1.py},
\nolinkurl{203_b1k3_zeta.py} with its CUE and per-prime checks
\nolinkurl{203_cue4.py}, \nolinkurl{203_b1k3_cue.py},
\nolinkurl{203_yerel4.py}); predictions committed before comparison
(\texttt{182190f}, \texttt{bf214f0}, \texttt{a056eb0}, \texttt{e5f2580}),
post-audit correction \texttt{62b36dd}, the test against zeros in
\nolinkurl{sifir_sinamasi_b1k3/}, Figure~\ref{M-fig:ladder}:
\nolinkurl{figures/ladder_ratios.py}; report
\nolinkurl{203_ORAN_KAPPA_RAPOR.md}.
\end{itemize}
The LMFDB zero files are not redistributed; their sha256 stamps are in
\texttt{ONKAYIT\_200C.json}.

\begin{thebibliography}{99}

\bibitem{companion1} U.~Sezen, \emph{The joint law of zero gaps and local
maxima of Hardy's $Z$-function}, preprint (2026); archived at Zenodo,
\href{https://doi.org/10.5281/zenodo.22942475}{doi:10.5281/zenodo.22942475}.

\bibitem{companion2} U.~Sezen, \emph{The prime-wave anatomy of the
gap--amplitude law for Riemann zeros}, preprint (2026); archived at Zenodo,
\href{https://doi.org/10.5281/zenodo.22942475}{doi:10.5281/zenodo.22942475}.

\bibitem{companion3} U.~Sezen, \emph{A response theory for the Riemann
zero gas}, preprint (2026); archived at Zenodo,
\href{https://doi.org/10.5281/zenodo.22942475}{doi:10.5281/zenodo.22942475}.

\bibitem{companion7} U.~Sezen, \emph{Arithmetic satellites of the Bragg
comb of the Riemann zero lattice: residue classes, characters, and a blind
test of the Bogomolny--Keating mirror law}, preprint (2026); archived at
Zenodo, \href{https://doi.org/10.5281/zenodo.22942475}{doi:10.5281/zenodo.22942475}.

\bibitem{Bar00} E.~W.~Barnes, \emph{The theory of the $G$-function},
Quart.\ J.\ Pure Appl.\ Math.\ \textbf{31} (1900), 264--314.

\bibitem{CG85} J.~B.~Conrey, A.~Ghosh, \emph{A mean value theorem for the
Riemann zeta-function at its relative extrema on the critical line},
J.\ London Math.\ Soc.\ (2) \textbf{32} (1985), 193--202.

\bibitem{Gonek93} S.~M.~Gonek, \emph{An explicit formula of Landau and its
applications to the theory of the zeta-function}, in: M.~Knopp et al.\
(eds.), A Tribute to Emil Grosswald: Number Theory and Related Analysis,
Contemp.\ Math.\ \textbf{143}, Amer.\ Math.\ Soc., Providence, RI, 1993,
395--413.

\bibitem{HLPC24} C.~Hughes, S.~Lugmayer, A.~Pearce-Crump, \emph{The second
moment of the Riemann zeta function at its local extrema}, J.\ London
Math.\ Soc.\ (2) \textbf{112} (2025), e70250, arXiv:2411.05573.

\bibitem{KS00} J.~P.~Keating, N.~C.~Snaith, \emph{Random matrix theory and
$\zeta(1/2+it)$}, Comm.\ Math.\ Phys.\ \textbf{214} (2000), 57--89.

\bibitem{Landau12} E.~Landau, \emph{\"Uber die Nullstellen der
Zetafunktion}, Math.\ Ann.\ \textbf{71} (1912), 548--564,
doi:\nolinkurl{10.1007/BF01456808}.

\bibitem{LMFDB} The LMFDB Collaboration, \emph{The $L$-functions and modular
forms database}, \url{https://www.lmfdb.org}, 2026, [Online; accessed
27 September 2026].

\bibitem{Mez07} F.~Mezzadri, \emph{How to generate random matrices from the
classical compact groups}, Notices Amer.\ Math.\ Soc.\ \textbf{54} (2007),
592--604, arXiv:math-ph/0609050.

\bibitem{Odl87} A.~M.~Odlyzko, \emph{On the distribution of spacings
between zeros of the zeta function}, Math.\ Comp.\ \textbf{48} (1987),
273--308.

\bibitem{OdlyzkoTables} A.~M.~Odlyzko, \emph{Tables of zeros of the Riemann
zeta function}, \url{https://www-users.cse.umn.edu/~odlyzko/zeta_tables/}.

\bibitem{Pla15} D.~J.~Platt, \emph{Computing $\pi(x)$ analytically},
Math.\ Comp.\ \textbf{84} (2015), no.~293, 1521--1535,
doi:\nolinkurl{10.1090/S0025-5718-2014-02884-6}.

\bibitem{Win12} B.~Winn, \emph{Derivative moments for characteristic
polynomials from the CUE}, Comm.\ Math.\ Phys.\ \textbf{315} (2012),
no.~2, 531--562, doi:\nolinkurl{10.1007/s00220-012-1512-1}; arXiv:1109.0227.

\end{thebibliography}

\end{document}
